
[Is there any semantic analysis of science papers?]
There are two types of analysis that can be done on scientific publications, semantic and relational.

\subsection{Citation Analysis} % Relational Analysis?
This is a study that finds patterns in different citation links to get insights such as mapping influence.\\*
These come in the form of 1) Link-Structure Rankings, 2) Pattern-Based Similarity, 3) Community \& Role Structure, 4) Temporal \& Evolutionary Citation Analytics, 5) Field-Normalisation \& Comparability, 6) Author/Entity-Level Indices, and 7) Counting \& Data Hygiene. A brief dive into each of them before a reflection from a GFSP perspective is given.

\subsubsection{Link-Structure Rankings}
\subsubsection{Pattern-Based Similarity}
\subsubsection{Community \& Role Structure}
\subsubsection{Temporal \& Evolutionary Citation Analytics}
\subsubsection{Field-Normalisation \& Comparability}
\subsubsection{Author/Entity-Level Indices}
\subsubsection{Counting \& Data Hygiene }


\subsubsection{Reflection}
These tend to miss semantic context, which is important as this holds the actual concepts, ideas, methods, and knowledge. They also hold no meaning or weightings in individual citations which induces a lot of noise that is only removed with large samples samples.


\subsection{Semantic Analysis}
[Intro words to semantic analysis]
This will explore: 1) Representation Foundations, 2) Semantic Similarity \& Mapping, 3) Concept \& Entity Normalisation, 4) Citation-Context Semantics, 5) Temporal \& Evolutionary Semantics, 6) Text-Graph Hybrids, and 7) Evaluation, Data, \& Hygiene.

\subsubsection{Representation Foundations}
\subsubsection{Semantic Similarity \& Mapping}
\subsubsection{Concept \& Entity Normalisation}
\subsubsection{Citation-Context Semantics}
\subsubsection{Temporal \& Evolutionary Semantics}
\subsubsection{Text-Graph Hybrids}
\subsubsection{Evaluation, Data, \& Hygiene}

\subsection{Methodological \& Theoretical Analysis}


\subsection{Traditional Metrics}
Citations, h-index, impact factor.

\textbf{Pros}: simple, widely used.

\textbf{Cons}: popularity bias, prestige effects, salami-slicing.

\subsection{Network-based Metrics}
PageRank, influence measures.

Better than counts, but still blind to semantic novelty.

Altmetrics \& Social Indicators: Twitter, Mendeley, etc.

Capture visibility, but amplify hype rather than intellectual content.

\subsection{Semantic Scientometrics}
Embeddings (topic models, SciBERT, SBERT, cosine similarity).

Start of capturing meaning, but usually measure similarity not inheritance/novelty.

\subsection{Paradigm Shift Detection}
Kuhn’s theory, attempts to quantify scientific revolutions.

Good framing, but not operationalised at semantic level.

\subsection{Gap Statement}
None of the above frameworks separates inherited vs novel contribution → motivation for your genetic evolutionary model.
